\chapter{Introduction}
\label{chap:introduction}

\section{Background and Motivation}

Urban bus planning has become more connected as electric buses enter fleets
that were once operated mainly by conventional vehicles.  For an internal
combustion engine (ICE) bus, the next trip is mostly a question of time and
reachability.  For an electric vehicle (EV), the same decision also depends on
remaining autonomy, possible charging time, and the availability of a free
charger \cite{li2014,kooten2017,rogge2018}.  A feasible timetable is therefore
not enough on its own.

These links are the focus of this thesis in the MINOA Senior challenge.  The
problem combines timetable selection, vehicle chaining, fleet assignment,
charging decisions, and terminal capacity, and these decisions influence each
other throughout the day.  A timetable choice changes the trips that must be
served; a vehicle connection changes the later battery state; and a charging
decision can occupy a limited terminal resource.  One challenge in this work
was that timetable feasibility and vehicle feasibility first looked separate,
but in the benchmark they were strongly linked.  The aim is therefore to
construct feasible and cost-efficient schedules while keeping these interactions
visible.

\section{Public Transport Planning Context}

Public transport planning is often divided into levels. Network and line
planning define the operated lines and terminals. Timetabling chooses the
passenger trips and their departure times. Vehicle scheduling then connects the
selected trips into vehicle blocks \cite{ceder2016,desaulniers2007,bunte2009}.
In a purely sequential process, the
timetable is fixed before the vehicle schedule is built. This is easier to
handle, but it can remove timetable alternatives that would need fewer vehicles
or create better charging opportunities \cite{michaelis2009,schmid2015}.

The MINOA setting sits at the boundary between timetabling and vehicle
scheduling. The input contains many candidate trips, but only a feasible subset
must be selected. The selected trips must satisfy headway rules, and every
selected trip must then be covered exactly once by a vehicle block. Because EV
blocks depend on autonomy and charging opportunities, a good solution must
coordinate the passenger-service side and the vehicle-operation side.
This is a small but important shift from a purely sequential workflow: the
timetable is not only a service plan, but also the input to the later vehicle
and charging decisions. This was not immediately obvious at the beginning,
because timetable feasibility and vehicle feasibility look like separate
questions until the larger instances are actually constructed.

\begin{figure}[H]
\centering
\includegraphics[width=\textwidth]{figures/fig37_research_positioning.png}
\caption{Planning story of the MINOA path-cover method.}
\label{fig:research-positioning}
\end{figure}
\FloatBarrier

\Cref{fig:research-positioning} summarizes the planning logic used throughout
the thesis. The solver first chooses a feasible timetable from the candidate
trip pool. It then asks a vehicle-scheduling question: which selected trip can
follow another selected trip? This question becomes a directed compatibility
graph. A path cover of this graph gives vehicle blocks. Finally, the solver
checks whether each block can be operated by an EV or must remain ICE, including
charging time, battery autonomy, and terminal capacity.  The main message is
that no single layer is sufficient on its own: a timetable must be operable, a
vehicle block must respect time and space, and an EV assignment must remain
feasible over the full block.  The same three decisions return later as
timetable paths, compatibility graphs, and mixed-fleet vehicle schedules.

\section{MINOA Challenge Setting}

The MINOA Senior input contains line directions, candidate trips, time windows,
terminal nodes, deadhead arcs, vehicle types, charging rates, cost
coefficients, and capacity limits. The output must contain the selected
timetable and complete vehicle blocks. Each block starts and ends at the depot,
contains selected passenger trips, and may include pull movements, waiting,
breaks, parking, and charging activities.

The computational study reports Small, Medium, and Large as the headline instances. It also
uses all available Senior instances as a wider benchmark. The additional
instances have different numbers of lines, directions, candidate trips, and
headway windows. They help to check whether the implementation works only for
the three headline files or also for a broader set of Senior inputs.

\section{Research Objectives and Questions}

The main objective is to design and evaluate a solver for integrated
timetabling and mixed-fleet vehicle scheduling. The solver should combine
timetable selection, vehicle chaining, fleet assignment, charging decisions,
and independent feasibility checking in one reproducible workflow.

The work is guided by three research questions:

\begin{enumerate}
    \item How can a feasible timetable be selected from the candidate trips
    while respecting headway rules and required boundary trips?
    \item How can selected trips be transformed into vehicle blocks in a way
    that reduces vehicles and operational cost?
    \item How much solution quality is gained when a simple greedy construction
    is extended to graph-based path cover, weighted path cover, and finally a
    multi-start path-cover matheuristic?
\end{enumerate}

\section{Scope and Delimitations}

The work focuses on the MINOA Senior problem data and its timetable,
vehicle-scheduling, charging, capacity, and cost rules. Crew scheduling,
stochastic delays, real-time dispatching, passenger assignment, and multi-depot
extensions are outside the scope. This restriction keeps the methodology close
to the given challenge and makes the computational results easier to interpret.

The method is heuristic for the full integrated problem and does not prove
global optimality. However, it uses exact graph subproblems inside the
construction, especially matching-based path cover
\cite{dulmage1958,edmonds1965,galil1986}. The reported schedules are feasible
solutions produced by a reproducible matheuristic. They can also provide useful
upper bounds and warm starts for a future exact method, as discussed in
\Cref{sec:future-global-optimality}.

\section{Research Contributions}

The main contributions are:

\begin{enumerate}
    \item It formulates the MINOA planning task using consistent notation for
    timetable selection, compatibility graphs, path-cover vehicle blocks,
    fleet assignment, battery propagation, charging, capacity, and cost.
    \item It implements four algorithmic variants: a greedy constructive
    baseline, an unweighted path-cover algorithm, a weighted path-cover
    algorithm, and a multi-start path-cover matheuristic. The final variant
    also uses bounded charging-aware edge scoring, EV charging look-ahead, and
    adaptive candidate-family settings.
    \item It gives detailed benchmark tables and figures for Small, Medium,
    Large, and all twelve Senior instances, including graph statistics, cost,
    vehicle count, EV share, selected-trip percentage, and runtime.
    \item It provides a reproducible implementation with a single experiment
    command and automated result reporting.
\end{enumerate}

\clearpage
\section{Structure of the Thesis}

The thesis is organized so that the reader first sees the planning context and
then reaches the algorithmic and computational parts.  This order matters
because the MINOA problem combines public-transport planning, graph-based
vehicle scheduling, and electric-bus feasibility.  \Cref{chap:literature}
starts with the research background.  It reviews timetabling, vehicle
scheduling, integrated planning, electric-bus scheduling, mixed fleets, and
graph-based optimization.  The focus is on the ideas needed for this thesis:
headway-based timetable selection, trip-to-trip compatibility, path covers,
matching, and charging-aware scheduling.

After this background, \Cref{chap:problem} introduces the MINOA Senior
benchmark in operational terms.  It describes the input data, candidate trips,
time windows, timetable rules, vehicle-block requirements, depot movements,
electric-vehicle autonomy restrictions, charging capacities, and cost
components.  \Cref{chap:formulation} then translates these planning rules into
mathematical notation.  It defines the compatibility graph, path-cover
structure, vehicle-type assignment, battery propagation, charging feasibility,
capacity restrictions, paid-break accounting, and timetable-variant scoring.
These definitions give the precise language used by the solution methodology.

The algorithmic part begins in \Cref{chap:methodology}.  The chapter develops
the four implemented approaches in increasing order of strength: greedy
construction, unweighted path cover, weighted path cover, and the final
multi-start path-cover matheuristic.  This sequence is useful because it shows
why the final method was chosen.  The greedy method gives a simple feasible
baseline, the path-cover methods use the compatibility graph more globally,
and the multi-start method adds timetable and edge-scoring variation.
\Cref{chap:implementation} connects these ideas to the Python software.  It
describes the repository structure, modules, experiment commands, hardware
setting, validation workflow, and reproducibility process.

The empirical part is presented in \Cref{chap:results}.  It reports the
validated results for Small, Medium, Large, and the full Senior benchmark,
including objective values, vehicle counts, EV/ICE fleet mix, runtime, lower
bounds, charging, and break minutes.  \Cref{chap:discussion} interprets the
evidence by comparing the benefits and limits of the methods.  Finally,
\Cref{chap:conclusion} summarizes the methodological contribution, the main
findings, and possible exact-certification extensions.
